\documentclass[12pt,a4paper]{article}

% ------------------ PAQUETES ------------------
\usepackage[utf8]{inputenc}
\usepackage[T1]{fontenc}
\usepackage[spanish,es-tabla]{babel} 
\usepackage{geometry}
\usepackage{setspace}
\usepackage{hyperref}
\usepackage{pgfgantt}
\usepackage{float}
\usepackage{enumitem} 
\usepackage{colortbl} % Para colorear celdas
\usepackage{xcolor}   
\usepackage{array}    % Nuevo paquete para mejorar un poco las tablas

\geometry{margin=2.5cm}
\setstretch{1.2}      
\definecolor{headergray}{gray}{0.9}

% ------------------ TÍTULO ------------------
\title{
    \vspace{-2cm}
    \Large \textbf{Planificación y Cronograma del TFG}\\[10pt]
    \large Diseño y análisis de un agente LLM autónomo para auditorías de ciberseguridad
}

\author{\textbf{Alejandro Barrena Millán}}
\date{17 de noviembre de 2025}

\begin{document}

\maketitle
\thispagestyle{empty} 
\hrule
\vspace{1em}

\section{Introducción}

El presente documento detalla la planificación temporal y la metodología de trabajo establecida para el desarrollo del Trabajo de Fin de Grado (TFG) titulado provisionalmente \textit{Diseño y análisis de un agente LLM autónomo para auditorías de ciberseguridad en entornos aislados}.

El objetivo es definir un calendario estructurado que garantice el cumplimiento de los hitos de desarrollo y la entrega de la memoria final en los plazos previstos.

% ============================================================
\section{Metodología: Scrum Adaptado}

El desarrollo del proyecto se regirá por un enfoque ágil basado en \textit{Scrum}, adaptado a las necesidades de un trabajo individual. Este marco de trabajo permite iterar sobre el proyecto, asegurar entregables funcionales en cada etapa y adaptar la planificación ante posibles imprevistos técnicos.

El ciclo de vida se divide en \textbf{Sprints} (ciclos de trabajo) con los siguientes eventos:
\begin{itemize}
    \item \textbf{Sprint Planning:} Definición de tareas y alcance al inicio de cada fase.
    \item \textbf{Desarrollo e Implementación:} Ejecución técnica de las tareas.
    \item \textbf{Sprint Review \& Retrospective:} Validación del entregable y ajuste de la metodología para el siguiente ciclo.
\end{itemize}

% ============================================================
\section{Planificación de Sprints y Tareas}

A continuación, se describe el desglose de tareas organizado por Sprints. 

\begin{description}[style=nextline, leftmargin=1.5cm, font=\bfseries]

    \item[Sprint 1: Definición y Organización (Semana 1)]
    Establecimiento de las bases del proyecto.
    \begin{itemize}[leftmargin=0.5cm, label=$\circ$]
        \item Definición detallada de requisitos funcionales y no funcionales.
        \item Selección y configuración de herramientas de gestión (e.g., Jira/Trello, GitHub).
        \item Investigación del estado del arte preliminar.
    \end{itemize}

    \item[Sprint 2: Entorno de Pruebas (Semanas 2 - 3)]
    Configuración de la infraestructura segura.
    \begin{itemize}[leftmargin=0.5cm, label=$\circ$]
        \item Despliegue del entorno de terminal en sandbox (Docker/VM).
        \item Desarrollo de la API de interfaz para la comunicación con el entorno.
        \item Pruebas de conectividad y seguridad del entorno aislado.
    \end{itemize}

    \item[Sprint 3: Integración del LLM (Semanas 4 - 6)]
    Conexión del motor de inteligencia artificial con el entorno.
    \begin{itemize}[leftmargin=0.5cm, label=$\circ$]
        \item Selección y configuración del modelo LLM (local o API externa).
        \item Desarrollo del módulo de interpretación de comandos.
        \item Prototipado inicial: el agente puede enviar comandos simples y recibir *outputs*.
    \end{itemize}

    \item[Sprint 4: Desarrollo del MVP (Semanas 7 - 8)]
    Implementación de la lógica de auditoría.
    \begin{itemize}[leftmargin=0.5cm, label=$\circ$]
        \item Programación de agentes autónomos con capacidad de decisión.
        \item Implementación de flujos de auditoría básica (e.g., escaneo de puertos, identificación de servicios).
        \item Validación del Producto Mínimo Viable (MVP).
    \end{itemize}

    \item[Sprint 5: Validación y Pruebas (Semana 9)]
    Consolidación del sistema.
    \begin{itemize}[leftmargin=0.5cm, label=$\circ$]
        \item Ejecución de auditorías *end-to-end* automatizadas.
        \item Análisis de falsos positivos y refinamiento de *prompts*.
        \item Pruebas de estrés y estabilidad.
    \end{itemize}

    \item[Sprint 6: Documentación y Memoria (Semanas 10 - 11)]
    Redacción académica.
    \begin{itemize}[leftmargin=0.5cm, label=$\circ$]
        \item Redacción de la memoria del TFG.
        \item Generación de gráficos, tablas de resultados y bibliografía.
        \item Revisión de formato y estilo.
    \end{itemize}

    \item[Sprint 7: Defensa (Semana 12)]
    Preparación final.
    \begin{itemize}[leftmargin=0.5cm, label=$\circ$]
        \item Diseño de diapositivas de presentación.
        \item Ensayo del discurso y preparación de la demo (si aplica).
    \end{itemize}

\end{description}

% ============================================================
\section{Cronograma Visual y Calendario}

Para facilitar el seguimiento del proyecto, se establece la correspondencia entre las semanas de trabajo y las fechas de calendario real, seguido de la representación visual mediante diagrama de Gantt.

\subsection{Calendario Detallado}

\begin{table}[H]
\centering
\renewcommand{\arraystretch}{1.3} % Más espacio entre filas
\begin{tabular}{|c|l|l|}
\hline
\rowcolor{headergray} \textbf{Semana} & \textbf{Fecha de Inicio (Lunes)} & \textbf{Fase / Sprint Activo} \\ \hline
1  & 17 de Noviembre, 2025 & S1: Definición y Organización \\ \hline
2  & 24 de Noviembre, 2025 & S2: Entorno Sandbox (Inicio) \\ \hline
3  & 01 de Diciembre, 2025 & S2: Entorno Sandbox (Cierre) \\ \hline
4  & 08 de Diciembre, 2025 & S3: Integración LLM (Inicio) \\ \hline
5  & 15 de Diciembre, 2025 & S3: Integración LLM (Desarrollo) \\ \hline
6  & 22 de Diciembre, 2025 & S3: Integración LLM (Cierre) \\ \hline
7  & 29 de Diciembre, 2025 & S4: Desarrollo MVP (Inicio) \\ \hline
8  & 05 de Enero, 2026     & S4: Desarrollo MVP (Cierre) \\ \hline
9  & 12 de Enero, 2026     & S5: Validación y Pruebas \\ \hline
10 & 19 de Enero, 2026     & S6: Documentación (Borrador) \\ \hline
11 & 26 de Enero, 2026     & S6: Documentación (Revisión) \\ \hline
12 & 02 de Febrero, 2026   & S7: Defensa y Presentación \\ \hline
\end{tabular}
\caption{Correspondencia de semanas con fechas naturales.}
\end{table}

\subsection{Diagrama de Gantt}

El diagrama ilustra la distribución temporal relativa de los Sprints a lo largo de las 12 semanas.

\begin{figure}[H]
\centering
\begin{ganttchart}[
    hgrid,
    vgrid,
    x unit=1.1cm,              
    y unit chart=0.6cm,
    bar/.append style={fill=blue!40, draw=none}, 
    bar height=0.6,
    group right shift=0,
    group top shift=0.7,
    group height=.3,
    title label font=\bfseries\footnotesize,
    bar label font=\footnotesize
]{1}{12}

% Títulos de las semanas
\gantttitle{Semanas de Desarrollo}{12} \\
\gantttitlelist{1,...,12}{1} \\

% Barras del cronograma
\ganttbar{S1: Organización}{1}{1} \\
\ganttbar{S2: Entorno Sandbox}{2}{3} \\
\ganttbar{S3: Integración LLM}{4}{6} \\ 
\ganttbar{S4: MVP Agente}{7}{8} \\
\ganttbar{S5: Validación}{9}{9} \\
\ganttbar{S6: Memoria}{10}{11} \\
\ganttbar{S7: Defensa}{12}{12}      

% Relaciones 
\ganttlink{elem0}{elem1}
\ganttlink{elem1}{elem2}
\ganttlink{elem2}{elem3}
\ganttlink{elem3}{elem4}
\ganttlink{elem4}{elem5}
\ganttlink{elem5}{elem6}

\end{ganttchart}
\caption{Cronograma general del proyecto.}
\end{figure}

% ============================================================
\section{Conclusión}

La planificación presentada asegura un flujo de trabajo constante y controlado. La estructura modular de los Sprints permite afrontar la complejidad de integrar Modelos de Lenguaje (LLM) en entornos de ciberseguridad paso a paso, mitigando riesgos técnicos antes de la fase final de documentación.

\end{document}